% Behavior Programs as Measures over Proofs — draft 0.5
% Build: tectonic metis.tex   (or pdflatex + bibtex-free; references are inline)
\documentclass[11pt]{article}

\usepackage[a4paper,margin=1.1in]{geometry}
\usepackage{amsmath,amssymb,amsthm}
\usepackage{mathtools}
\usepackage{stmaryrd}
\usepackage{microtype}
\usepackage{booktabs}
\usepackage{enumitem}
\usepackage{xcolor}
\definecolor{linkblue}{rgb}{0.1,0.15,0.45}
\usepackage[colorlinks=true,linkcolor=linkblue,citecolor=linkblue,urlcolor=linkblue]{hyperref}
\usepackage{cmap}

% ---- notation -------------------------------------------------------------
\newcommand{\llp}{\textsc{llp}}
\newcommand{\metisc}{\texttt{metisc}}
\newcommand{\lolli}{\multimap}
\newcommand{\tensor}{\otimes}
\newcommand{\bang}{\mathord{!}}
\newcommand{\parr}{\bindnasrepma}    % ⅋ (reversed ampersand), from stmaryrd
\newcommand{\sem}[1]{\llbracket #1 \rrbracket}
\newcommand{\ff}[2]{#1^{\underline{#2}}}          % falling factorial
\newcommand{\R}{\mathbb{R}}
\newcommand{\N}{\mathbb{N}}
\newcommand{\Rpos}{\R_{+}}
\newcommand{\mllb}{\mathrm{MLL}^{\mathrm B}}
\newcommand{\cclass}{\mathcal{F}}                  % the factorizable class
\newcommand{\repourl}{\url{https://github.com/metis-lang-dev/metis-lang}}

% status markers
\newcommand{\stproved}{[\textsf{proved}]}
\newcommand{\stmech}{[\textsf{mechanized}]}
\newcommand{\stpinned}{[\textsf{pinned}]}
\newcommand{\stsketch}{[\textsf{sketch}]}
\newcommand{\stopen}{[\textsf{open}]}

\theoremstyle{plain}
\newtheorem{theorem}{Theorem}
\newtheorem{proposition}[theorem]{Proposition}
\newtheorem{conjecture}[theorem]{Conjecture}
\newtheorem{corollary}[theorem]{Corollary}
\newtheorem{lemma}{Lemma}
\theoremstyle{definition}
\newtheorem*{definition}{Definition}
\theoremstyle{remark}
\newtheorem*{remark}{Remark}

\title{\textbf{Behavior Programs as Measures over Proofs}\\[1ex]
\large Staged stochastic multiset rewriting with a linear-logic front end,
an $\mllb$ proof-net intermediate representation, and an executable,
certified, machine-checked compilation pipeline}
\author{The metis project\thanks{Sources, corpus, goldens, and the Lean
mechanization: \repourl. Correspondence: \texttt{info@robofortis.com}.}}
\date{Draft 0.5 --- 2026-10-29}

\begin{document}
\maketitle

\begingroup\small\noindent
\textbf{Status markers.} \stproved{} = proof in this paper;
\stmech{} = machine-checked in Lean~4 in the public repository
(directory \texttt{formal/}, \repourl); \stpinned{} = mechanically
test-pinned in the repository; \stsketch{} = proof sketch, mechanization
open; \stopen{} = open problem.

\medskip\noindent
\textbf{Changes from 0.3.} A new section on belief
(\S\ref{sec:belief}): epistemic probability enters a program only
through a finite port, a belief is a mixture over initial multisets
(Theorem~\ref{thm:mixture}), the program reads it as a level atom, and
its certificate gains a finite sheet and loses nothing
(Proposition~\ref{prop:cells}, Corollary~\ref{cor:cert}). Interfaces,
packages and compositions are now part of the language
(\S\ref{sec:packages}), with the composite measure a product that the
tick protocol determines (Proposition~\ref{prop:composite}). The
evaluation bound separates grounding, which is a join, from inference,
which is the proof-net's (\S\ref{sec:evaluation}). The text has been
rewritten for a reader seeking understanding: the history of the
drafts and the engineering particulars are in the repository's record,
not here.
\endgroup

\begin{abstract}
We present a behavior-planning framework in which a \emph{program} --- a
catalog of typed, staged, weighted multiset-rewriting rules in the Ceptre
lineage of intuitionistic linear logic --- denotes a \emph{probability
measure over its own derivations}. Three readouts of one object coexist:
proof theory (which derivations exist), operational semantics (how they
interleave), and measure (how likely each is). We prove that the
interpreter's run distribution and the denoted $k$-step dynamic Bayesian
network agree exactly (run-level preservation), characterize the image of
the denotation as stage-gated mass-action stochastic Petri nets
(equivalence onto the image, now machine-checked in Lean~4), and show
that for a statically-recognizable \emph{factorizable class} the measure
descends to the Mazurkiewicz trace quotient --- runs weigh their
occurrence nets, not their interleavings --- connecting our compiled
representation to the Bayesian proof-nets of Ehrhard, Faggian and
Di~Guardia read in the opposite direction: they evaluate one net where we
select among many. Uncertainty about the world --- epistemic, as
opposed to the aleatoric weights inside a program --- enters only
through a finite port: a belief is a mixture over initial multisets, the
program reads it as a level atom, and its certificate gains a finite
sheet and loses nothing. The correspondence is not only semantic but
\emph{compiled}: programs lower through a typed ground normal form to a
proof-net-shaped intermediate representation (\textsc{choice} sites as
CPT boxes, SSA atom occurrences as wires) on which typing --- including
machine-checked well-formedness as EFD-style Bayesian proof-nets --- cost
certificates and canonical hashing are IR-level analyses, and from which
exact filtering, junction-tree elimination, fixed-budget sampling, tree
search, a dependency-free C++ interpreter, and standard factor-graph
export are generated back-ends. Admission-time cost certificates make the
pipeline deployable in real-time physical systems: programs are rejected
or downgraded before evaluation, never killed mid-flight. We report exact
parity chains (interpreter $\doteq$ filter $\doteq$ factored $\doteq$
exported artifact $\doteq$ C++ interpreter, at $10^{-9}$ across 200+
gates), provable optimality on IPPC benchmark instances, a growing Lean~4
mechanization of the metatheory, and honest limits (the $2^n$ support
cliff and what the factored evaluator does and does not buy back).
\end{abstract}

\section{Introduction}\label{sec:intro}

Two research programs meet in the middle of this paper.

Ehrhard, Faggian and Di~Guardia (EFD) read a correspondence between
Bayesian networks and multiplicative linear logic \textbf{right to left}:
a BN's joint distribution \emph{is} an $\mllb$ proof-net (MLL plus
probabilistic box axioms carrying CPTs), and exact inference \emph{is}
the inductive evaluation of that net in probabilistic coherence spaces,
cost-matched to factor elimination \cite{EFD24,EFD26}. One net, given by
the model; the computational question is evaluation.

We read the same correspondence \textbf{left to right}. A linear-logic
program does not denote one proof; it \emph{generates a space of
derivations} --- plans, explanations, futures. Equipping the
\textsc{choice} points of the operational semantics with weights makes
that space a probability space, and the interesting computations are over
the ensemble: \emph{sum} over derivations is filtering and prediction;
\emph{argmax} is most-probable explanation and planning;
\emph{reweighting by evidence} is diagnosis and behavior identification.
Probability selects among proofs; proof theory generates them.

\paragraph{Two kinds of uncertainty.} The weights on \textsc{choice}
are the dice inside a behavior: given the world, this is how the
program acts. Call that part \emph{aleatoric}; it is fixed when the
program is certified, and tuned, when it is tuned, by counts. A second
kind of uncertainty is about the world itself --- where the cup is,
whether the beast is hostile --- and it is \emph{epistemic}: it belongs
to the agent's map and changes with every observation. A program with
fixed weights is a certified Markov decision process. With the map's
belief entering through a port, it is a belief MDP --- over a
\emph{finite} quantisation, because the belief enters as a finite atom
(\S\ref{sec:belief}). Keeping the two kinds apart, and letting the
second cross into the first only through a finite port, is what lets
every certificate of this paper survive the arrival of belief.

\paragraph{What the lollipop buys.} For the experienced reader it
pays to separate, at the door, the two halves of the logic--probability
correspondence, because they are statements of different kinds. The
\emph{Boolean} half --- events form an algebra, probability is a
measure on it, $\wedge/\vee/\neg$ are $\cap/\cup/\complement$ --- is
classical \cite{Kol33}, geometrically exact, and structurally inert:
it says which events exist and yields a flat joint measure, but
nothing about how the joint \emph{factors}. Everything substantive in
this paper lives in the other half: the linear implication
$A \lolli B$ read as the conditional $P(B \mid A)$. This is not
another Venn-style equivalence; it is a statement about
\emph{direction and factorization}, which the event algebra cannot
express at all. A BN node with its parents is one lollipop --- EFD's
box axiom: negative premises, one positive conclusion, a CPT, with
$A \lolli B = A^\perp \parr B$ \cite{Gir87} the reason implication
and box share a shape; the chain rule
$P(x_1 \ldots x_n) = \prod_i P(x_i \mid \mathrm{pa}_i)$ is a
composition of lollipops (in Markov-categorical terms: conditionals
are directed morphisms and composition is the chain rule
\cite{CJ19,Fri20}); a DBN's transition kernel is one lollipop per
slice. The reading is load-bearing three times over in what follows.
The emission compiles every clause $c_r \lolli o_r$ into a CPT row
whose parents are exactly the clause's premises, and the box-rule
shape is machine-checked on every artifact
(Proposition~\ref{prop:bpn}). The run measure factorizes per site
into conditionals given exactly those premises
(Theorem~\ref{thm:trace}) --- and that factorization, not the flat
measure, is what the cost certificates, junction trees and exports
operate on. And the implication is \emph{consuming}: co-enabled
clauses compete for the same tokens, so conditioning on an outcome
anticorrelates the alternatives --- explaining-away \emph{is}
resource competition, priced by the falling factorials of
Theorem~\ref{thm:preservation} --- a phenomenon the Boolean half is
structurally blind to, since classical contraction lets a premise be
used twice. The direction of $\lolli$ is likewise why conditioning
remains a semantic operation on the measure rather than a syntactic
one on proofs (\S\ref{sec:obstructions}, O3).

The two directions meet at the occurrence net: each of our runs,
quotiented by interleaving, is an EFD-style net; our $k$-step measure is
a weighted superposition of such nets; conditioning and MAP pick nets out
of the superposition --- a question that cannot arise on the EFD side,
where a query has exactly one net.

\paragraph{The mission, stated plainly.} The project's goal is a
\emph{formally provable} program --- in both senses --- from linear logic
to probability graphs: a language whose denotation into factored
probabilistic models is established by theorems, whose theorems are
progressively machine-checked in a proof assistant, and whose compiled
artifacts are certified by admission-time checkers so that what the
deployed interpreter evaluates is exactly what the theorems govern.
The theorems are on paper, the artifact checks run at every build, and
the metatheory is progressively discharged in Lean~4
(\S\ref{sec:mech}). The three levels of assurance --- paper proof,
mechanized proof, mechanically pinned artifact --- are marked
per-statement throughout.

This paper contributes, on the theory side:
\begin{enumerate}[leftmargin=2em]
\item \textbf{Run-level preservation} (Theorem~\ref{thm:preservation},
  \stproved{}, counting kernel \stmech{}): the interpreter's run
  distribution and the denoted $k$-step DBN are related by a
  measure-preserving bijection; the falling-factorial \textsc{choice}
  weights are exactly the token-individuation counts of the
  ordered-resource interpreter.
\item \textbf{Exactness of reachable-support filtering}
  (Theorem~\ref{thm:filter}, \stproved{}, \stmech{}): forward filtering
  over the reachable support is exact inference on the denoted DBN, not
  an approximation.
\item \textbf{Equivalence onto the image} (Theorem~\ref{thm:image},
  \stproved{}, \stmech{}): the denotation is an equivalence between
  ground staged programs modulo structural congruence and finite
  stage-gated mass-action stochastic Petri nets --- the honest
  ``isomorphism theorem''; logic $\leftrightarrow$ all-of-PGM is
  obstructed and we say exactly why (\S\ref{sec:obstructions}).
\item \textbf{Trace invariance on the factorizable class}
  (Theorem~\ref{thm:trace}, \stproved{}, descent skeleton \stmech{}):
  for programs admitted by our factored emission, the run weight
  factorizes per \textsc{choice} site and is invariant under
  Mazurkiewicz trace equivalence --- the measure is a measure on
  occurrence nets. Outside the class, per-slice normalizers are
  interleaving-dependent and the quotient fails in general; the class
  boundary is \emph{statically recognizable and enforced}.
\item \textbf{Quiescence is compiled away}
  (Proposition~\ref{prop:quiescence}, \stproved{}): the stage/quiescence
  discipline --- the one construct of our source language outside MLL
  --- is eliminated by the finite-horizon normal form for admitted
  programs: the IR is quiescence-free. The non-MLL residue is a
  \emph{source-language} convenience, not a semantic commitment.
  Moreover the emitted IR is a well-typed \textbf{Bayesian proof-net} in
  the EFD sense, machine-checked at every serialization
  (Proposition~\ref{prop:bpn}, \stproved{}, \stpinned{}) --- the static
  well-formedness half of the Di~Guardia typing question, whose
  compositional half remains Conjecture~\ref{conj:typing}.
\item \textbf{Belief through a finite port} (Theorem~\ref{thm:mixture},
  \stproved{}, \stmech{}; Proposition~\ref{prop:cells},
  Corollary~\ref{cor:cert}, \stproved{}):
  a categorical belief bound to a port makes the one-tick number a
  mixture over initial multisets; with the belief entering through the
  presence draw alone the number is multilinear, so a finite cell sheet
  bounds it; the program's certificate is unchanged.
\item \textbf{The composite measure is a product}
  (Proposition~\ref{prop:composite}, \stproved{}): under the tick
  protocol a composition of certified packages has the product of their
  measures, so nothing composite need be grounded to certify it.
\end{enumerate}

And on the systems side:
\begin{enumerate}[leftmargin=2em]\setcounter{enumi}{7}
\item A \textbf{compilation pipeline} (\S\ref{sec:pipeline}) with the
  shape of knowledge compilation: \texttt{.llp} source $\to$ typed
  ground/unrolled normal form $\to$ $\mllb$-shaped IR (sites + SSA
  occurrences) $\to$ generated back-ends. Slogan: \emph{$\mllb$ is to
  linear-logic multiset rewriting what SDD circuits are to probabilistic
  logic programming.}
\item \textbf{Admission-time cost certificates} (\S\ref{sec:certs}):
  induced width and total factor-ops computed symbolically from IR
  scopes before any table exists --- the real-time defense doctrine
  (reject or downgrade the evaluator ladder at admission; never kill
  mid-flight).
\item A \textbf{verification-translator architecture}
  (\S\ref{sec:cpp}): the compiler (\metisc, OCaml) and the Python
  referee toolkit are certifier and cross-check; the deployed C++
  runtime is a $\sim$300-line dependency-free interpreter of the
  compiled artifact, pinned by golden parity and by external referees.
\item A \textbf{mechanization} (\S\ref{sec:mech}): a Lean~4 development,
  shipped in the public language repository, machine-checking
  Theorem~\ref{thm:image} in full and the cores of
  Theorems~\ref{thm:preservation}, \ref{thm:filter}
  and~\ref{thm:trace}. The paper contains no Lean code; statements
  cited as \stmech{} name their Lean identifiers in \S\ref{sec:mech}.
\end{enumerate}

What is \emph{not} claimed: novelty of the static
BN$\leftrightarrow$proof-net correspondence (that is EFD's, published
first --- see \S\ref{sec:related} for the audit trail); an isomorphism
between the logic and all of PGM (obstructed,
\S\ref{sec:obstructions}); completeness of the learning loop (doctrine
and seams exist; two helpers remain unbuilt, \S\ref{sec:learning}).

\section{The source language and its operational semantics}
\label{sec:language}

\subsection{Programs}

A \textbf{catalog} declares finite types, namespaces with layer rights,
predicate signatures, backward (Horn) predicates, ports, and entries.
Entries are clause \emph{schemas} and link \emph{schemas} over typed
variables, each carrying a mandatory human-approvable comment and a
provenance record (\texttt{hand} $\mid$ \texttt{surprise} $\mid$
\texttt{template}). \textbf{Grounding} $\Pi$-instantiates schemas over
the signature, prunes by Horn-derivable guards and distinctness
constraints, and yields a \textbf{ground program}
\[
  P = (S,\ (C_s)_{s\in S},\ L,\ s_0,\ m_0)
\]
with stages $S$ (ordered), ground clauses $C_s$ per stage, ground links
$L$, initial stage $s_0$ and initial multiset $m_0$. A ground clause
\[
  r:\quad c_r \lolli o_r \quad [\text{persist } p_r] \quad
  \text{weight } w_r
\]
consumes multiset $c_r$, reads (bang-style, never decremented) $p_r$,
produces $o_r$, and carries a prior weight $w_r \in \Rpos$, optionally
bound through a \textbf{weight port}: a named factor whose implementation
lives in a registry, with a \emph{declared read-scope}
\texttt{reads \{\dots\}} (\S\ref{sec:emission}). A link is a clause that
additionally switches the active stage and is gated on quiescence.

\subsection{States and the three-tier CHOICE}\label{sec:choice}

A state is $\sigma = (s, m)$: active stage and multiset. One step:
\begin{itemize}[leftmargin=2em]
\item \textbf{Tier 0 (clauses).} Enabled clauses of $s$ draw
  \[
    W_r(m) = w_r \cdot \varphi(m, r), \qquad
    \varphi(m, r) = \prod_a \ff{m(a)}{c_r(a)} \cdot \prod_a m(a)^{p_r(a)}
  \]
  the falling-factorial transition count (consume) times the power count
  (reads). One enabled clause is chosen with probability
  $W_r / \sum_{r'} W_{r'}$. A \textbf{policy hook} may reweight or prune
  this tier (safety shields, behavior rows); the measure renormalizes.
\item \textbf{Tier 1 (links).} If no clause of $s$ is enabled
  (\emph{quiescence}), enabled links of $s$ compete the same way and
  switch the stage.
\item \textbf{Tier 2 (done).} Otherwise the state is absorbing.
\end{itemize}
\textsc{update} is deterministic given the chosen event. A \textbf{run}
of horizon $K$ is the sequence of $K$ chosen events (with \texttt{done}
self-loops); the \textbf{run measure} $\mu_K$ is the product of its
steps' \textsc{choice} probabilities.

This is the Ceptre operational core \cite{Mar15} --- itself the
first-order Horn fragment of intuitionistic linear logic as multiset
rewriting \cite{CS09} --- extended with (i) the weight priors as
first-class syntax, (ii) weight ports with declared scopes, and (iii) the
stage/quiescence discipline as tiered \textsc{choice}.

\subsection{Three readouts, one object}

The same $P$ supports: \textbf{proof theory} --- runs are derivations;
each run's dependency structure is an occurrence net
(\S\ref{sec:trace}); \textbf{operations} --- the interpreter samples
runs; \textbf{measure} --- $\mu_K$, with filtering, smoothing, MAP,
hypothesis posteriors and surprise as queries. Nothing is layered on top
of anything: all three are readouts of $P$.

\subsection{Interfaces, packages, compositions}\label{sec:packages}

How is a large behavior written? Not as one program. A robot's
navigation, grasping and task logic are separate \emph{specialists},
each a program of its own, and a \emph{director} that hands goals down
and reads reports up. The language keeps them separate and fixes, in
declarations, exactly what crosses between them.

\paragraph{The forward declaration is small.} An \textbf{interface
package} is a lifted file with only declarations: abstract types and
predicates over them, no rules, no facts, no stages, with its own name
and canonical key. Every predicate two programs share is declared
\emph{once}, there, and carries a \emph{nature}: a \emph{signal} (a
world reading or a crossing vocabulary --- $\mathit{goto}$,
$\mathit{reached}$), a \emph{resource} (the hands, the position: held
by the director, granted as a lease), or a \emph{report} (status, done,
stuck). A specialist includes the interface packages it speaks and
states its \emph{direction} on each predicate --- in, out, resource, or
world-private --- and a static check refuses a direction the nature does
not allow, a shared predicate declared locally, or two interfaces
declaring one name. The interface packages every composition speaks
--- the tick protocol, the director tree --- are ordinary lifted files
with no special status in the language; they are the language's
\emph{standard units}, shipped with each toolchain and resolved by the
include rule like any package (the including directory, the project
root, the toolchain's own units, then the user's shared packages), and
the link records which file each name resolved to.

\paragraph{The package is the certified unit.} A \textbf{package} is a
program with its manifest, its cases (an initial multiset, a script of
deliveries, the expected atoms) and its goldens. Its certificate is the
program's canonical key plus its cases passing against one
\emph{generic} interface stub that knows, from the manifest alone, what
to feed (the in-set), what to collect (the out-set), what to grant and
how to score. A specialist's certificate therefore never depends on the
director it will meet: the protocol it speaks is a declaration package
with no rules.

\paragraph{The shared atoms are the runtime's.} A \textbf{composition}
is a set of packages over interface packages under a \emph{tick
protocol} with \emph{copy semantics}: every port a member reads is a
copy taken at the start of that member's own tick, what a member emits
at a tick is visible to any reader --- itself included --- no earlier
than the reader's next tick, the relay between members is deterministic,
and each member samples its own stream. Members tick at their own
clocks (navigation fast, the person's intents slow), so there is no
beat shared across them to latch at; the rule is stated per reader, and
it is the same rule whether the holder of a port is in the reader's
process, in another process, or on another machine. Shared state is
never a live shared truth: a linear-logic $\bang p$ readable by every
rule and consumed by none would make a trigger that never settles and a
slow member that sees a moving target. The composition's certificate is its package key, the verdict
of the interface checks, and each member's own key in manifest order
--- nothing composite is grounded, because the composite measure is a
product the protocol determines (Proposition~\ref{prop:composite}).

\paragraph{The image is the composition's deployable denotation.} A
\textbf{bundle} is one package compiled: its ground wire, its sidecar,
its tool declarations, under a content hash. An \textbf{image} is a
composition \emph{linked}: every member's bundle, every shared
predicate resolved to its producers, consumers and nature, the capacity
and the priority order checked, and the director \emph{generated} from
the composition's declarations rather than written by hand, all under
one hash over a canonical file list. The link is where the static side
of the tick protocol is decided: that the members are exactly the
declared specialists, that each world predicate consumed is one the
world provides, that no belief crosses a member boundary, that the
response the director promises is reachable in ticks, that one planner
exists. What the image does \emph{not} contain is a \textbf{plan}: the
assignment of members and world ports to processes, and the carrier
between them (sockets today; a robot middleware later). A plan is
checked against the image, is linked by the image's hash, and stays
outside it, because under copy semantics a plan changes only
\emph{delay} --- a copy that crossed a process boundary arrives a tick
later, within the response budget the link checked --- and never what
any member computes from the copies it holds. The one thing a split
does not preserve is a global order across the boundary, which is why a
property stated across processes is stated against the message order
and not against a beat.

This is what certification attaches to. A certificate names an image
hash and nothing else; a runtime reads the image and nothing else ---
not the composition, not the manifests, not the source. A member may be
replaced iff it was certified against the same interface keys and the
composition links again, a static and cheap test; a host library
upgrade changes a member's binding in the image and not the
composition's keys. The same image loads into one process holding every
member or into two processes with the person's dialog on one side and
the behavior on the other, and the measured difference on a two-cook
kitchen was the delay alone: every member's fired rules identical beat
for beat under the two plans, the only rows differing being the two
first-beat fetches whose command crossed the boundary a tick later.

\section{Run-level preservation}\label{sec:preservation}

The denoted \textbf{$k$-step DBN} has variables $X@t$ (the event chosen
at step $t$, including \texttt{quiesce}/\texttt{done}) and $K@t$ (the
multiset state), with conditionals \emph{defined} by
\S\ref{sec:choice}'s \textsc{choice} and deterministic \textsc{update}.
The substantive theorem relates this to the \emph{finer-grained}
semantics an ordered-resource interpreter (Ceptre's) actually executes,
where a step chooses not a ground clause but an \textbf{instance}: a
clause together with an injective assignment of concrete resource tokens
to its consume slots and a choice of token per read occurrence.

\begin{theorem}[preservation; \stproved{}, counting kernel \stmech{}]
\label{thm:preservation}
Fix a horizon $K$. The map that sends an interpreter trajectory
(sequence of instances) to its event sequence (sequence of ground
clauses) pushes the interpreter's trajectory distribution forward to
$\mu_K$, and the induced map on state trajectories is a
measure-preserving bijection between the interpreter process and the DBN
process.
\end{theorem}

\begin{proof}
Both processes share state space and deterministic \textsc{update}, so
it suffices to show per-step agreement of the choice distributions; the
trajectory statement follows by induction on $K$ with the product
decomposition of both measures.

Fix $\sigma = (s, m)$ at tier 0 (tiers 1--2 are identical on both sides,
as links carry the same count structure and \texttt{done} is
deterministic). On the interpreter side, tokens are individuated: the
multiset $m$ is a set of tokens with $m(a)$ tokens of species $a$. An
instance of clause $r$ chooses, for each species $a$, an \emph{injective}
assignment of $c_r(a)$ consume slots to distinct tokens of $a$ --- there
are $m(a)\,(m(a)-1)\cdots(m(a)-c_r(a)+1) = \ff{m(a)}{c_r(a)}$ such
assignments --- and, independently, one token per read occurrence ---
$m(a)^{p_r(a)}$ choices. (Equivalently: an \emph{unordered} choice of
$c_r(a)$ distinct tokens followed by an ordering onto the slots,
$\binom{m(a)}{c_r(a)} \cdot c_r(a)! = \ff{m(a)}{c_r(a)}$ --- the count
is the same whether one individuates slots or only tokens, which is why
the falling factorial is forced and not a convention.) The instance
count of $r$ at $m$ is therefore exactly $\varphi(m, r)$. The
interpreter draws instances with probability proportional to the clause
prior $w_r$ per instance; hence the total probability of \emph{some}
instance of $r$ is proportional to $w_r \cdot \varphi(m, r) = W_r(m)$,
which is the DBN's \textsc{choice} weight. All instances of a fixed $r$
produce the same successor multiset (removing $c_r$, adding $o_r$ is
species-level), so the pushforward along instance $\to$ clause is exact,
and distinct clauses map to their own $W_r$ mass. Normalizing over the
enabled set gives equality of the one-step kernels.
\end{proof}

Two corollaries worth naming. First, \emph{substitution collapse
pricing}: when grounding identifies two variables ($X{=}Y$), the ground
clause's consume multiset gains multiplicity 2 and $\varphi$ prices it
as a falling factorial --- the theorem shows this is forced by token
individuation, not a modeling choice. Second, \emph{explaining away =
resource competition}: the collider at $X@t$ (parents = the consumed
species) anti-correlates alternative consumers of the same tokens under
conditioning, priced exactly by the same counts. \stpinned{}: oracle
parity $\le 10^{-12}$ on the analytic fixtures, $\le 10^{-9}$ on the
road corpus, across all horizons. The counting kernel --- instances of
$r$ at $m$ are equinumerous with $\varphi(m,r)$, and the
instance-to-clause pushforward carries weight $w_r\,\varphi(m,r)$ ---
is \stmech{} (\S\ref{sec:mech}).

\begin{theorem}[exact filtering; \stproved{}, \stmech{}]
\label{thm:filter}
Let $D_t$ be the forward filter's distribution over states after $t$
steps, computed over the reachable support only. Then $D_t$ equals the
DBN's exact marginal at slice $t$.
\end{theorem}

\begin{proof}
Induction on $t$. $D_0$ is the point mass at $(s_0, m_0)$. The DBN's
slice-$(t{+}1)$ marginal is $\sum_\sigma P(\sigma' \mid \sigma)\,
D_t(\sigma)$ by the law of total probability; because \textsc{update} is
deterministic, $P(\sigma' \mid \sigma)$ is supported on finitely many
$\sigma'$ per $\sigma$, and every $\sigma$ with $D_t(\sigma) > 0$ is by
construction in the filter's support. The filter computes exactly this
sum. No mass ever leaves the reachable set, so restriction loses
nothing.
\end{proof}

The point of stating the triviality: the filter is \emph{not} an
approximation of PGM inference --- it is variable elimination in the one
elimination order (time-forward) whose intermediate factors are
supported on reachable states. Its failure mode is honest: support grows
like $2^n$ on synchronous-parallel domains (the IPPC cliff,
\S\ref{sec:cliff}), which is what the factored back-end addresses.

\section{Canonical identity: programs and runs}\label{sec:canonical}

Two identities, two quotients \stpinned{}:

\paragraph{Programs mod presentation.} \texttt{canonical()}
content-sorts clauses, links and names; \texttt{program\_key()} hashes
the result (stage order and initial multiset retained --- they are
semantics). Invariant under presentation shuffles; discriminates weight,
multiplicity, stage wiring, init and dropped-clause mutations. This key
is the cache key for every compiled artifact (\S\ref{sec:certs}) ---
canonical identity is not bookkeeping, it is what makes AOT compilation
sound.

\paragraph{Runs mod interleaving.} Independence of events is footprint
disjointness --- writes (consume $\cup$ produce) conflict with
everything, reads conflict only with writes --- with the stage
discipline encoded as an implicit resource (clauses read
\texttt{stage@s}; links write both endpoint stage tokens), so
cross-stage ordering falls out of the same rule. \texttt{run\_key()}
hashes the \textbf{Foata normal form} of the trace: the canonical
layering of the run's occurrence net. Two interleavings of the same
concurrent run share one key; symmetric ground instances collapse --- a
quotient we want.

The delicate question --- \emph{does the measure respect this
quotient?} --- is \S\ref{sec:trace}, and the answer is one of the
paper's main theorems.

\section{Equivalence onto the image}\label{sec:image}

No isomorphism between the logic and all of PGM exists at the
implementable level; \S\ref{sec:obstructions} gives the three
obstructions. The honest strong statement characterizes the image.

\begin{definition}[the image class]
A \textbf{stage-gated mass-action SPN} is a tuple
$N = (A, \bar S, T, G, s_0, m_0)$ where $A$ is a finite set of places
(ground atoms); $\bar S$ is a finite \emph{ordered list} of stage labels
(stages with no transitions are retained --- they are declared
structure, not derived); $T$ is a finite \textbf{multiset} of
\emph{stage transitions}, each a quadruple $(s, c, p, o, w)$ of a stage
label, consume/read/produce arc multiplicity vectors over $A$, and a
rate constant $w \in \Rpos$ (a multiset, because nothing forbids two
indistinguishable transitions, and the firing rule counts them
separately); $G$ is a finite multiset of \emph{gate transitions}
$(s, s', c, p, o, w)$ additionally carrying a target stage;
$s_0 \in \bar S$ and $m_0$ is a marking. The propensity of a transition
at marking $m$ is $w$ times the mass-action count (falling factorials on
consumed places, powers on read places), and the firing rule is
\S\ref{sec:choice} verbatim: stage transitions of the active stage
compete first, gates fire at stage quiescence and set the register,
normalization per step.
\end{definition}

\begin{definition}[structural congruence]
$P \equiv P'$ iff $P'$ is obtained from $P$ by bijectively renaming
clause and link \emph{names} and reordering the clause and link lists.
Atoms, stage labels, stage order, weights, and the initial data are
fixed pointwise --- they are semantics, not presentation (renaming an
atom or a stage changes which SPN is denoted; renaming a clause does
not).
\end{definition}

\begin{theorem}[equivalence onto the image; \stproved{}, \stmech{}]
\label{thm:image}
The denotation $\sem{\cdot}$ from ground staged programs to stage-gated
mass-action SPNs is (i) total and functional, and constant on
congruence classes; (ii) surjective; (iii) injective up to structural
congruence. Hence it descends to a bijection between
programs-mod-congruence and the class.
\end{theorem}

\begin{proof}
(i) \emph{Totality and invariance.} $\sem{P}$ is the
\S\ref{sec:language} construction: places = the (finitely many) ground
atoms of the signature, $\bar S$ = the declared stage list in declared
order, each ground clause $r$ of stage $s$ contributes the stage
transition $(s, c_r, p_r, o_r, w_r)$, each link the corresponding gate,
$s_0$ and $m_0$ carried over. Every component of the tuple is assembled
from the \emph{content} of clauses and links --- names appear nowhere,
and $T$, $G$ are multisets, so list order appears nowhere either. Hence
$\sem{\cdot}$ is total, functional, and constant on congruence classes.

(ii) \emph{Surjectivity.} Given $N$ in the class, define a reading-back
$U(N)$: fix any enumeration $t_1, \dots, t_k$ of the multiset $T$
(duplicates enumerated with multiplicity) and any enumeration
$g_1, \dots, g_l$ of $G$; let $U(N)$ have stages $\bar S$ in order,
clause $r_i$ named freshly (say $c_i$) with stage, multisets and weight
copied from $t_i$, link $l_j$ named freshly with endpoints and data
copied from $g_j$, and initial data $(s_0, m_0)$. Then
$\sem{U(N)} = N$ \emph{on the nose}: each tuple component round-trips
by construction, and the operational behaviors agree because the
class's firing rule was \emph{defined} to be \S\ref{sec:choice}. (The
choice of enumeration changes $U(N)$ only up to $\equiv$, which (i)
already quotients away.)

(iii) \emph{Injectivity up to congruence.} Suppose
$\sem{P} = \sem{P'} = N$. From $N$ we recover: the stage list with its
order (so $P$ and $P'$ declare the same stages, including empty ones),
$s_0$ and $m_0$, and --- per stage label $s$ --- the multiset of
content quadruples $(c, p, o, w)$ of the clauses of $s$, and the
multiset of gate contents $(s, s', c, p, o, w)$ of the links. For each
stage $s$, the clause lists of $P$ and $P'$ over $s$ induce equal
content multisets, so there is a bijection between them matching equal
content; where several clauses share identical content, any matching of
the duplicates will do, because identical-content clauses are fully
interchangeable --- swapping them is itself a congruence. The same
argument gives a per-endpoint-pair matching of links: a link's
endpoints $(s, s')$ are part of its \emph{content}, so the matching
automatically respects the stage graph --- this is the point where
renaming might have interacted with the stage-linking structure, and it
cannot: clause and link names occur nowhere else in a program (clauses
do not reference each other or links by name; links reference
\emph{stages}, and stages are not renamed), so the per-stage and
per-link matchings assemble into one global renaming-plus-reordering
with no cross-boundary constraint left to check. That renaming exhibits
$P \equiv P'$.
\end{proof}

Two boundary remarks the proof depends on. \emph{Empty stages}: $\bar S$
retains them; had the SPN carried stages only as the partition induced
by $T$, a declared-but-unpopulated stage would be forgotten and (iii)
would fail (two programs differing by an unreachable empty stage would
collide) --- the class definition is chosen so that the denotation
forgets \emph{exactly} presentation. \emph{Duplicate transitions}: $T$
is a multiset for the same reason --- two identical-content clauses
double the propensity of that content, which a set-valued $T$ would
silently halve. The theorem is \stmech{} (\S\ref{sec:mech}): programs as
name-carrying lists, the class as multiset-valued tuples, congruence as
renaming plus permutation, with totality, invariance, surjectivity (an
explicit fresh-naming section of $\sem{\cdot}$) and
injectivity-up-to-congruence each a named Lean result. The \stpinned{}
evidence is a fuzz test of the same statement on the implementation:
the canonical key is invariant under presentation shuffles and
discriminates every semantic mutation.

Classical context: the $\tensor$-fragment $\leftrightarrow$ Petri
reachability correspondence is Gunter--Gehlot \cite{GG89} and
Engberg--Winskel \cite{EW93}; MSR as the Horn fragment of ILL is
Cervesato--Scedrov \cite{CS09}; the mass-action propensity shape is
Gillespie's \cite{Gil77}. The assembly --- \emph{staged}, \emph{gated},
\emph{discrete-time normalized}, with an executable preservation test
and a mechanized equivalence --- is ours.

\section{The measure on occurrence nets: trace invariance}
\label{sec:trace}

Here is the problem the seed analysis left open. The engine weights
$W_r(m)$ are local, but each step's \emph{normalizer} $\sum W_{r'}$
ranges over everything enabled at that step --- and what is co-enabled
depends on the interleaving. In general, two interleavings of the same
occurrence net can carry different products of normalized
probabilities: the measure lives on sequences, not nets.
(Abbes--Benveniste's branching cells \cite{AB06} are the general
repair; importing them wholesale is \stopen{}.)

The factored-emission work yields a large, statically-recognizable
class where the repair is free.

\begin{definition}[factorizable class $\cclass$]
A ground program (with its horizon and declared port scopes) is in
$\cclass$ iff the factored emission (\S\ref{sec:emission}) admits it:
(a) all multiset counts are 0/1; (b) the stage/tier schedule is
branch-independent (which micro-step runs which competition is the same
on every run); (c) at each micro-step, the candidate set forms a
\textbf{site}: every possibly-enabled event consumes the same anchor
token, and each candidate's enabledness and weight depend only on a
\emph{scope} --- its consumed/read atoms plus its port's declared
\texttt{reads} scope; (d) no event produces an atom that may already be
present (checked statically or by the runtime mass check).
\end{definition}

\begin{theorem}[trace invariance on $\cclass$; \stproved{}, descent
skeleton \stmech{}]\label{thm:trace}
For $P \in \cclass$, the weight of a run factorizes as
\[
  \mu(\mathit{run}) \;=\; \prod_{\text{sites } \sigma \text{ visited}}
    P(\mathit{chosen}_\sigma \mid \text{scope-valuation}_\sigma)
\]
and is invariant under Mazurkiewicz trace equivalence. Consequently
$\mu_K$ descends to a measure on occurrence nets: \texttt{run\_key}
classes are measurable atoms, and the $k$-step measure is a weighted sum
over EFD-style nets.
\end{theorem}

\begin{proof}
By (b), every run visits the same sequence of sites, so the product
ranges over a common index set; it remains to show each factor is
well-defined on the trace class. Fix a site $\sigma$ and consider two
runs in one trace class. Their prefixes before $\sigma$ differ by a
sequence of adjacent transpositions of independent events
(Mazurkiewicz).

\emph{Lemma (footprint stability).} If events $e, e'$ are independent
(disjoint footprints in the \S\ref{sec:canonical} sense, including
stage tokens), then executing $e'$ does not change (i) the truth of any
atom in $e$'s site scope, nor (ii) $e$'s enabledness. --- Immediate:
$e$'s scope atoms are among its consumed/read/declared atoms;
independence says $e'$ writes none of them.

By the lemma and induction on transpositions, the scope valuation at
site $\sigma$ is equal in both runs. The site's normalizer is
$\sum_{\text{candidates}} (\text{enabledness} \times w \times
\text{port-factor})$, and by (c) every term is a function of that scope
valuation alone --- \emph{this is exactly what the declared
\texttt{reads} contract buys}: an opaque binding reading undeclared
global state would break the lemma, which is why the emission rejects
undeclared ports. Hence each factor agrees; the products agree. With
(a), $\varphi \equiv 1$, so no falling-factorial term reintroduces
marking dependence.
\end{proof}

Three remarks. First, the theorem \emph{identifies our sites with
branching cells} on $\cclass$: local competition arenas whose choices
are independent of scheduling --- ``concurrency = probabilistic
independence'' made constructive by a compiler check. Second, the class
boundary is \emph{enforced, not assumed}: static admission plus the
runtime total-mass check ($Z = 1$) convert any reachable violation into
an error rather than a silent approximation (\stpinned{}: rejection
gates). Third, outside $\cclass$ the engine semantics remains
well-defined (Theorems~\ref{thm:preservation}--\ref{thm:filter}); what
fails is only the descent to nets --- normalizers become
schedule-dependent, which is precisely the interleaving junk the
general theory must quotient by force. The descent skeleton --- run
weights that factor through schedule-stable local factors are invariant
under adjacent transposition of independent events, hence constant on
Mazurkiewicz classes --- is \stmech{} in abstract form
(\S\ref{sec:mech}); the site/scope instantiation for the concrete
emission remains on paper.


\paragraph{The composite measure factorizes by the tick.} The same
descent holds one level up, for a composition of packages
(\S\ref{sec:packages}).

\begin{proposition}[the composite measure is a product; \stproved{}]
\label{prop:composite}
Let a composition run under the tick protocol: (F1) within tick $t$ a
beating member $d$ reads only its private state and the copies taken for
it at boundary $t$; (F2) nothing $d$ emits at $t$ is visible to another
member before boundary $t{+}1$; (F3) the copy, relay and grant transfer
at a boundary is a deterministic function of the members' latched
outputs and the world's readings; (F4) the members' random streams are
disjoint. Then for every composite trace $\tau$,
\[
  P(\tau) \;=\; \prod_t \ \prod_{d \text{ beats at } t}
  P_d\big(\text{steps of } d \text{ at } t \;\big|\; \mathit{state}_d(t)\big),
\]
where $\mathit{state}_d(t)$ is $d$'s private state with its boundary-$t$
copies, fixed by $\tau$'s prefix through (F3), and each $P_d$ is the
member's own run measure $\mu_K$ of \S\ref{sec:choice}.
\end{proposition}

\begin{proof}
By (F1) and (F2), no member's \textsc{choice} within a tick depends on
another member's \textsc{choice} in the same tick; by (F4) the draws are
independent; by (F3) the state each member starts a tick from is a
function of the prefix. The probability of the tick's joint step is
therefore the product over beating members of their own one-tick
measures, and the product over ticks follows by the chain rule.
\end{proof}

What the product does \emph{not} give is the flat program's measure:
interleavings within a tick are not represented --- (F2) forbids them
--- and a chain of three members takes two boundaries and shows as such.
That is the price of certifying members alone, and it is paid in
latency, not in semantics. Arbitration is the director's
\textsc{choice} at its own beat, inside its own measure; the grant it
emits reaches the holder at the next boundary.

\section{Why no logic \texorpdfstring{$\leftrightarrow$}{<->} PGM
isomorphism, and what quiescence really is}\label{sec:obstructions}

\textbf{O1 --- interleaving junk.} Sequence-level artifacts can only be
homomorphic-with-preservation; the isomorphism-grade statement lives at
the proof-net quotient --- exactly where EFD obtain theirs, and where
Theorem~\ref{thm:trace} puts us for $\cclass$.

\textbf{O2 --- copy/discard asymmetry.} The PGM world is a Markov
category \cite{Fri20,CJ19}: copying and marginalization are free.
Linear logic licenses contraction/weakening only under $\bang$ --- our
persistent read is the licensed copy; dropped products are weakening. A
functor from a resource-sensitive syntax into a semicartesian target
cannot be inverted on the whole target; the honest repair is
Theorem~\ref{thm:image}'s image characterization.

\textbf{O3 --- inference direction.} $\lolli$ fixes the generative
direction; a BN's defining extra power is conditioning (Bayesian
inversion). Our system computes posteriors \emph{about} derivations
(smoothing, MAP, hypothesis posteriors, all exact over the reachable
support \stpinned{}) but the lollipop has no proof-theoretic inverse
here --- inversion is a semantic operation on the measure, not a
syntactic one on proofs.

\paragraph{Quiescence beyond MLL --- and
Proposition~\ref{prop:quiescence}.} The stage-token encoding
(\S\ref{sec:canonical}) shows stages alone are MLL-typeable sugar:
clauses read \texttt{stage@s}, links consume/produce stage tokens. The
residue is the quiescence gate --- ``fire only if \emph{nothing} in
this stage can'' --- a negative side-condition (inhibitor arc;
negation-as-failure) invisible to any multiset of premises, hence
outside MLL and outside the EFD setting. This was our honest deviation.
The compiled view dissolves it:

\begin{proposition}[quiescence is compiled away on $\cclass$;
\stproved{}]\label{prop:quiescence}
For $P \in \cclass$ at horizon $K$, tier decisions are
branch-independent, so the normal form resolves every quiescence test
statically; the emitted IR (sites + SSA occurrences + CPT boxes)
contains no quiescence construct, and is $\mllb$-shaped: boxes = sites
(negative premises = scope occurrences, positive conclusion = outcome),
contraction = SSA fan-out of an occurrence read by several later
scopes, weakening = occurrences never read again.
\end{proposition}

\begin{proof}
By $\cclass$(b) the schedule --- including which micro-steps are
link-tier --- is a function of the step index alone, computed once by
the admission analysis; the emission then never consults enabledness at
evaluation time. \stpinned{} by the emission gates: mixed-tier programs
are rejected, admitted programs carry no tier logic in the artifact.
\end{proof}

So ``quiescence beyond MLL'' resolves into: a \emph{source-level}
convenience (genuinely non-MLL, kept because staged authoring is how
humans write behavior) that the compiler eliminates for the class it
certifies. The general-class proof theory of quiescence (fixpoints?
subexponentials? inhibitor semantics?) remains \stopen{}.

\begin{proposition}[the emitted nets are Bayesian proof-nets;
\stproved{}, \stpinned{}]\label{prop:bpn}
For every admitted program and horizon, the emitted IR, read as an
$\mllb$ typed graph --- \textsc{choice} sites and SSA writes as
probabilistic boxes (negative premises = the parent occurrences of the
declared scope, positive main conclusion = the outcome/occurrence
variable), \texttt{unknowns} lifts as parentless prior boxes,
double-production guards as observed constraint boxes, atom fan-out as
contraction on negative atoms, unread occurrences as the net's
conclusions $\Delta$ --- is a \textbf{Bayesian proof-net} in the sense
of \cite[Defs.~3.1--3.2]{EFD26}: (a) every box satisfies the box rule
(all conclusions atomic, exactly one positive, $X \notin \{Y_i\}$); (b)
box conclusions are pairwise distinct; (c) every premise edge is the
conclusion of some box (no pending premises); (d) the net is correct
--- by \cite[Lemma~3.1]{EFD26} it suffices that the polarized
orientation is a DAG; (e) all conclusions in $\Delta$ are positive.
Site CPT rows are normalized or zero (stall rows, convertible to errors
only by the runtime $Z = 1$ check, \S\ref{sec:emission}); write boxes
are deterministic one-hot; guards are 0/1 indicators.
\end{proposition}

\begin{proof}
Each condition holds by construction of the emission
(\S\ref{sec:emission}): (b) is SSA single-writer --- every occurrence
variable is minted exactly once, every site variable is fresh per
micro-step; (c) holds because a variable enters a scope only through
the environment map, whose non-constant entries are exactly the minted
variables; (a) holds because a site never reads its own outcome and a
write never reads the occurrence it mints; (d) holds because every
premise of a box minted at micro-step $t$ was minted strictly earlier,
so emission order is a topological order of the polarized orientation;
(e) holds because pending edges are unread box main conclusions, which
are positive by the box rule.
\end{proof}

The point of \emph{also} pinning it: the construction argument is about
the emission as designed; the checker verifies conditions (a)--(e)
mechanically on the emitted artifact itself --- admission-time,
symbolic, zero binding calls, same doctrine as the cost certificate ---
and runs both over the IR golden corpus (plus the unknowns and guard
constructions) and \emph{at every IR serialization}: the emitter
refuses to write an ill-typed artifact, so the deployed C++ interpreter
--- which by design never re-types (\S\ref{sec:cpp}) --- only ever
receives certified nets. Adversarial mutations (a self-parent box, a
double writer, a dangling premise, a polarized cycle, a sub-normalized
CPT row) are each rejected, so the gate is not vacuous. Any future
emission bug that breaks well-typedness fails CI and the wire, not a
reviewer.

\begin{conjecture}[Di~Guardia compositional typing]\label{conj:typing}
Proposition~\ref{prop:bpn} discharges the \emph{static} half of the
typing question: our emitted nets are well-formed Bayesian proof-nets,
machine-checked. What remains is the \emph{compositional} half ---
checking the emitted nets against the well-behaved-composition typing
derivations of \cite{DEF26} (their type system governs how nets compose
via cut; our sites are their boxes, SSA fan-out their contraction, and
Proposition~\ref{prop:quiescence} removes the structural moves the
dynamic quiescence gate would have needed, so we conjecture every
emitted net types). Mechanizing their typing rules over our artifact is
the concrete next theory task. \stopen{}
\end{conjecture}

\section{Belief: epistemic probability through a finite port}
\label{sec:belief}

What does a program do when the world is not fully known? Everything so
far prices \emph{what the program does}: the \textsc{choice} weights are
the aleatoric part, the dice inside the behavior. A second kind of
uncertainty lives outside the program --- \emph{what the world is} ---
and it is epistemic: it changes as the agent observes. This section
says how the second kind enters without touching the first, and why the
certificates survive its arrival.

\subsection{The example, and the two missing factors}

A robot must fetch a cup. The cup is in one of three receptacles, and
the agent's map gives a belief: $0.6$ in the drawer, $0.3$ on the shelf,
$0.1$ in the bin. The program has a rule that sends the robot to a
receptacle,
\[
  \mathit{go}:\ \ \mathit{want}(O) \tensor \mathit{start} \tensor
  \bang\mathit{inrecep}(O,R) \lolli \mathit{want}(O) \tensor
  \mathit{goto}(R) \qquad \texttt{@w look}(O,R),
\]
whose weight is a \textbf{port}: a named factor the host binds, here to
the belief. With the map asserting $\mathit{inrecep}$ for every place
the cup might be, the rule is enabled three times and the port prices
the choice among them. This is the right program and the wrong
probability. A \textsc{choice} weight is normalized over the enabled
events of a state (\S\ref{sec:choice}); it answers \emph{where does the
robot go first}, not \emph{where is the cup}. Two things a total
probability needs are absent:
\begin{itemize}[leftmargin=2em]
\item \emph{The world factor.} ``Where is the cup'' is one distribution
  over three mutually exclusive facts. Three independent weights describe
  $2^3$ worlds, including the cup in two places and in none.
\item \emph{The product.} $P(\text{fetched within } h) = \sum_R
  P(\text{cup in } R)\cdot P(\text{fetched within } h \mid \text{cup in }
  R)$. The first factor is the belief; the second is the program's own
  measure. Nothing in \S\ref{sec:language}--\ref{sec:trace} multiplies
  them: a run starts from one initial multiset.
\end{itemize}
The rest of the section supplies the two factors (\S\ref{sec:catport},
\S\ref{sec:instrument}), brings the product back into the program as a
finite atom (\S\ref{sec:levels}), and proves that the program's
certificate is unchanged by the arrival (\S\ref{sec:beliefthm}).

\subsection{The categorical port}\label{sec:catport}

A \textbf{belief group} is declared over a port that already exists, by
naming which of the port's arguments is the one that varies:
\emph{for each object $O$, exactly one receptacle $R$ holds, and
$\mathtt{look}(O,\cdot)$ is its distribution}. The group's key is the
other arguments; its members are the values of the named one. Three
forms cover what the maps produce: \texttt{one} (exactly one member
holds; the masses sum to $1$), \texttt{lone} (at most one; the residual
is the world where none holds --- the cup is lost), and an independent
atom (present with its own probability). The host binds the port as
before, one value per argument tuple; the only new obligation is the
sum, and it is checked exactly, never silently renormalized. A
distribution the host did not mean is refused, not repaired.

The declaration lives beside the program, not in it. A belief is a fact
\emph{about the program's inputs}; the tick --- the one step the
interpreter samples --- never reads it, so the program and its canonical
key (\S\ref{sec:canonical}) do not move. What moves is the certificate's
key, which takes the canonical line of the belief declarations
(\S\ref{sec:beliefthm}).

In proof-net terms (\S\ref{sec:obstructions}, Proposition~\ref{prop:bpn})
a group is \emph{one} categorical box $\vdash X^{+}$ whose variable has
the group's $k$ members as values and the port as its CPT. Three
Bernoulli boxes denote the same distribution only once an
``exactly one'' factor couples all three; the categorical box is one
variable of cardinality $k$. The correct total and the cheaper net are
the same object.

\subsection{Belief facts and the world mixture}\label{sec:worlds}

In the instrument of \S\ref{sec:instrument} the presence of a belief
pred's atoms is not read from the initial multiset: it is \emph{drawn}
from the port. Each \texttt{one} group puts exactly one member present,
with its probability; each \texttt{lone} group one or none; each
independent atom is present or absent. A \textbf{world} $w$ is one
outcome of all the draws, with probability $P(w)$ the product over
groups, and the instrument starts from the mixture
\[
  \mu_0 \;=\; \sum_w P(w)\,\delta\big((s_0,\ m_0 \uplus w)\big).
\]
Two disciplines keep this honest. A case's initial multiset may not list
a belief atom (the draw puts it there). And belief preds are resource
preds, never Horn guards: a guard is eliminated at grounding
(\S\ref{sec:evaluation}), and a world must not change the ground
program --- every world shares one program and differs only in its
initial multiset.

A read of a belief fact is conditioned on the truth. In the world where
the cup is on the shelf, a rule enabled by $\bang\mathit{inrecep}(O,
\mathit{shelf})$ fires there and nowhere else. That is right for the
\emph{world's} rules (the observation ``absent'', the find). It is
omniscient for the \emph{agent's} decisions: a task whose $\mathit{go}$
reads the belief fact walks straight to the true place in every world,
and its success is overstated. The honest program keeps two preds ---
the agent's \emph{support} (where the cup may be, read by the decision)
and the belief fact (where it is, read by the world) --- and one port
prices both: as a \textsc{choice} weight on the decision and as the
presence of the truth. The instrument lists every rule that reads a
belief pred, so the author sees which decisions are conditioned on the
hidden truth.

\subsection{Level ports: the only belief-derived input of the tick}
\label{sec:levels}

The program never reads a number. A \textbf{level port} turns a number
--- a belief's value, or a query's answer --- into a finite atom by
declared thresholds: with cut points $0 < a_1 < \dots < a_n < 1$ and
names $\ell_0, \dots, \ell_n$, the value $p$ has level $\ell_i$ iff
$a_i \le p < a_{i+1}$ (with $a_0 = 0$, $a_{n+1} = \infty$), compared
exactly. The host asserts exactly one level atom per key and retracts
the previous one; it is a world reading like any other. A tick program
that reads only levels and static weights is a \emph{finite} program,
and finiteness is what \S\ref{sec:beliefthm} certifies.

This is the whole design in one sentence: the belief enters the
program as a finite atom, so the program with belief is a program of the
kind every theorem above already governs; the arithmetic that produced
the atom lives in an instrument beside the program, where it can be
exact.

\subsection{The one-tick instrument}\label{sec:instrument}

A \textbf{query} names an outcome and a horizon: the probability that
$Q$ holds at some step $t \le h$ (first passage; $Q$ read as wanted or
as a violation), or its presence at step $h$. The instrument answers for
the beliefs \emph{as they stand at this tick}:
\begin{itemize}[leftmargin=2em]
\item \emph{Beliefs are constant over the horizon.} The port table is
  evaluated once, before step $0$, and every weight port prices its
  events by that frozen table at every step. The map's own dynamics ---
  the cup carried away, the plant drying --- and its updates from what
  was seen stay in the map; the next tick's table comes from there.
\item \emph{Filtering stays in the map.} Within the horizon, what the
  program observes is conditioned exactly by the enumeration itself: in
  world $w$ every observation follows $w$'s truth.
\end{itemize}
The number is a sum-product over worlds and traces: products along a
trace (the draws, then the normalized \textsc{choice}s of
\S\ref{sec:choice}), sums where states meet; a state where $Q$ holds is
absorbed, its mass added to the answer. It is the forward filter of
Theorem~\ref{thm:filter} started from a mixture instead of a point.

The state budget is the cut. The filter retains at most $N$ states,
expanded in canonical order; the mass of what cannot be expanded is set
aside, and the answer is an interval
\[
  \mathit{lo} = \mathit{HIT}, \qquad
  \mathit{hi} = \mathit{HIT} + \mathit{CUT},
\]
sound for first passage because the cut mass may or may not reach $Q$
and nothing else can (\stmech{}: the dropped mass is moved by a
stochastic kernel, never created or destroyed). A query is answered in
bounded work at any horizon, and the width $\mathit{hi} - \mathit{lo}$ says how much the
budget hid. With no cut, $\mathit{lo} = \mathit{hi}$ is the exact
number --- an exact rational: every weight, port value and threshold
enters as the exact decimal of its shortest round-trip representation,
so $0.7$ is $7/10$ and a distribution written $0.7,\,0.3$ sums to one,
and the result does not depend on summation order. The brute-force
referee that enumerates every path agrees to the last digit by
construction.

\subsection{What is proved}\label{sec:beliefthm}

\begin{theorem}[a port is a mixture; \stproved{}, \stmech{}]\label{thm:mixture}
Let $P$ be a ground program with belief groups $B$ bound by $f$, and let
$\mu_w$ be the first-passage measure of $P$ started from
$(s_0, m_0 \uplus w)$ under the frozen table $f$. For every query $Q$,
horizon $h$ and budget $N$,
\[
  \mathit{lo} \;\le\; \sum_w P_f(w)\,\mu_w(Q \text{ within } h) \;\le\;
  \mathit{hi},
\]
with equality on both sides --- the instrument's number is the mixture
--- when no state is cut.
\end{theorem}

\begin{proof}
The instrument's recursion is Theorem~\ref{thm:filter}'s forward filter
applied to the distribution $\mu_0 = \sum_w P_f(w)\,\delta(s_0, m_0
\uplus w)$ with $Q$-states absorbing. The one-step kernel is the same
for every world (the worlds share the ground program and the frozen
table), so the filter is linear in its initial distribution:
$F^h(\sum_w P_f(w)\,\delta_w) = \sum_w P_f(w)\,F^h(\delta_w)$, and
absorption commutes with the sum because it is a per-state operation.
Each summand is $\mu_w(Q \text{ within } h)$ by
Theorem~\ref{thm:filter}. With a cut, the retained recursion is a
sub-filter that drops nonnegative mass; the dropped mass is exactly
$\mathit{CUT}$, and since the kernel is stochastic, what it would have
contributed to $\mathit{HIT}$ lies in $[0, \mathit{CUT}]$.
\end{proof}

The theorem says that a quantised port adds nothing the static model
could not already express: it names \emph{which} initial multisets a
belief chooses and with what mass. Everything the earlier sections
proved about one run from one initial multiset is inherited by the
mixture, term by term.

\begin{proposition}[multilinearity, and the cell table; \stproved{}]
\label{prop:cells}
Suppose the belief enters the number only through the presence draw
(no \textsc{choice} weight reads a belief port). Then
$\sum_w P_f(w)\,\mu_w(Q)$ is multilinear in the groups' distributions:
linear in each $f(g, \cdot)$ separately. Consequently, over any product
of polytopes --- one per group, cut out by the level thresholds of
\S\ref{sec:levels} --- its minimum and maximum are attained at products
of vertices.
\end{proposition}

\begin{proof}
$P_f(w) = \prod_g f(g, w_g)$ is a product with one factor per group, and
$\mu_w$ does not depend on $f$ by hypothesis. A function of the form
$\sum_w \prod_g f(g, w_g)\,c_w$ is linear in each factor's distribution
when the others are held fixed. A multilinear function on a product of
polytopes is, in each coordinate block, linear, so an extremum over the
product is attained at a point extremal in every block, i.e.\ at a
product of vertices.
\end{proof}

A \textbf{cell} is one level per level port. For each cell, the
certificate records the interval $[\mathit{lo}_c, \mathit{hi}_c]$ that
the instrument gives for each query over every belief consistent with
the cell: by Proposition~\ref{prop:cells} this is the minimum and maximum
over the cell's vertex worlds, each one exact instrument run. With no
belief level at all, the vertices are the \emph{pure} worlds and the
bound is $\min_w$ and $\max_w$ of $\mu_w(Q)$ --- the numbers the brute
referee already enumerates. The sheet is finite: a \texttt{one} group of
$k$ members confined to a box contributes at most $k \cdot 2^{k-1}$
vertices, an independent atom two.

\begin{corollary}[certification is preserved]\label{cor:cert}
The certificate of a program with belief ports is the program's own
certificate (\S\ref{sec:certs}, keyed by \texttt{program\_key}) together
with the finite cell sheet, keyed additionally by the canonical line of
the belief and level declarations. The port selects rows of a table
that was certified without it; it never adds a cell.
\end{corollary}

The honest boundary is where the hypothesis of
Proposition~\ref{prop:cells} fails: a \textsc{choice} weight that reads
a belief port (the fetch rule's \texttt{@w look}) makes the trace
measure itself depend on the belief as a \emph{ratio} of port values.
The function is then rational, not multilinear, and the vertex bound is
not sound. Such a cell is listed \emph{uncertified (choice reads a
belief)}. Both forms are legal; the certificate tells them apart, and
the certified form gates the decision on a belief \emph{level} with
static weights.

\subsection{The rewrites: on the instrument's net, never on the tick}

The tick has no net; it samples one step. The instrument's net is
static --- the beliefs, the ground program, $h$ and $Q$ are fixed before
it runs --- so the proof-net rewrites of EFD \cite{EFD26} apply to it,
each keeping the exact rational and each checked against the brute
referee: \emph{weakening} (a group no ground event touches sums to one
and is dropped), \emph{deferred draws} (a group is drawn at the first
state where its presence could matter, exact because the draw is
independent of the trace so far), \emph{absorption} (a state where $Q$
holds stops), \emph{merging} (equal states sum --- the frontier is the
cut-net's interface), and \emph{the cone} (restriction to the outcome's
influence cone when its closure holds). None of them is a semantic
choice; each is Theorem~\ref{thm:mixture}'s sum evaluated in a cheaper
order.

\subsection{Ports by step: the map's forecast as rows}\label{sec:bystep}

A belief that moves along the horizon --- a beast approaching, a car
changing lane --- is a hidden chain, and tracking it inside the
instrument would be filtering, which \S\ref{sec:instrument} leaves in
the map. The cheap form that pays today: the host computes its forecast
\emph{outside} the instrument, one row per step, and the instrument
reads row $t$ at step $t$. A rule factor prices the event chosen at step
$t$ by row $t$; a presence draw and a belief level read row $0$; a row
past the table's end repeats the last, so the frozen table is the
one-row case and nothing earlier changes meaning.

Two points decide what a row may carry. A rule factor is a \emph{per-step
hazard} --- the probability of the move at this step given where things
are --- never a marginal: a marginal (``nearer than now at step $t$'')
read as a factor at every step compounds. And a quantity constant over
the horizon --- an intent: hostile, wandering, fleeing --- is better a
categorical \emph{belief} drawn once per world than a by-step port: the
intent then persists along each trace (a hostile beast closes at every
step, not at the mixture's average rate), Proposition~\ref{prop:cells}
holds for it, and only the hazards conditioned on the intent move by
step:
\[
  \mathit{close}:\ \ \mathit{beast}(B,Q) \tensor \mathit{next}(P,Q)
  \tensor \bang\mathit{intent}(B,I) \lolli \mathit{beast}(B,P)
  \qquad \texttt{@w closing}(B,I).
\]
On a lane with a beast believed at three cells, the risk of meeting it
within eight events under a frozen hazard, a rising hazard by step, and
the last row held frozen are $0.207$, $0.294$ and $0.346$: the moving
belief lies between its first and last rows, as it must.

With a by-step port, multilinearity is lost: the rows multiply along
every trace and are functions of the forecast. The certificate then
holds \emph{over the initial levels} --- the cells are the tick's
levels, as before --- and the forecast that turns the tick's belief into
rows is part of the sealed image (\S\ref{sec:packages}), listed as
\emph{certified over the initial levels, predict in the image}, never
by the vertex bound.

\subsection{The aim: ports as functions, and the horizon}

The form this section stops short of is a port bound as a \emph{function}
of the step and the state --- the belief as a hidden Markov chain inside
the instrument, updated by what each trace observes. The exact rational
semantics and the cut carry over to it; Proposition~\ref{prop:cells}
does not, and the certificate would have to be stated over the
forecast's own parameters. The language refuses it today, by name, so
that every program that compiles is one the cell sheet covers. It is
open problem~5 (\S\ref{sec:open}).

\section{The compiled pipeline (systems)}\label{sec:pipeline}

\begin{center}\small
\begin{tabular}{l}
\texttt{.llp} source (rulescript: Ceptre-core rules, typed header,
ports,\\
\quad includes with RDDL-style domain/instance lifting)\\
\quad $\downarrow$ parse $\to$ lossless AST
($\texttt{parse}(\texttt{pretty}(a)) = a$)\\
\quad $\downarrow$ compile $\to$ catalog + ports; typecheck; admission\\
\qquad (packs: containment by construction --- a pack may only
\emph{add},\\
\qquad\ never touch base namespaces)\\
typed ground/unrolled \textbf{normal form} \quad $\leftarrow$
stages/quiescence lowered away (Prop.~\ref{prop:quiescence})\\
\quad $\downarrow$\\
$\mllb$-shaped IR: occurrence net = \textsc{choice} sites (categorical
CPT boxes\\
\quad over declared scopes) + SSA atom occurrences (binary wires)\\
\quad $\downarrow$ IR-level analyses: contextual typing; Bayesian
proof-net check (Prop.~\ref{prop:bpn});\\
\qquad canonical hash (\texttt{program\_key}); \textsc{cost
certificate} (induced width + factor-ops,\\
\qquad symbolic --- \S\ref{sec:certs}); junction tree (AOT)\\
\quad $\downarrow$ generated back-ends:\\
\qquad exact reachable-support filter (+ policy hook) \hfill
[Thm~\ref{thm:filter}]\\
\qquad factored elimination / junction tree (mid-scale exact)\\
\qquad fixed-$N$ sampler (deterministic cost, same measure)\\
\qquad THTS planner (same measure, budgeted)\\
\qquad C++ IR interpreter (deployment; \S\ref{sec:cpp})\\
\qquad BN/factor-graph export (UAI + DOT; \S\ref{sec:export})\\
\end{tabular}
\end{center}

\subsection{The knowledge-compilation frame}

The nearest cousin is ProbLog $\to$ d-DNNF/SDD: compile once into a
form whose evaluation is polynomial in the form's size, amortized over
queries. \textbf{$\mllb$ is to LL-MSR what SDD circuits are to PLP} ---
with the difference that our target form is also a \emph{published
semantics} (EFD's nets), not only a data structure. To our knowledge no
prior system compiles a staged stochastic multiset-rewriting language
into proof-net-shaped factor graphs; EFD's own pipeline starts from
BNs.

\subsection{Evaluation: where the cost is, and what bounds it}
\label{sec:evaluation}

Where does a program's evaluation spend its time? In two stages that
obey two different bounds, and it pays to keep them apart because only
one of them is covered by the proof-net theory.

\paragraph{Grounding is a join, before any net exists.} A rule $r$ has
typed variables $x_1 \ldots x_n$ and guards $g_1 \ldots g_m$ (Horn
premises). Its ground events are the substitutions for which every
guard is derivable:
\[
  |\mathrm{ground}(r)| \;=\; \big|\{\theta \in \textstyle\prod_i \tau_i
  : \forall j.\ \Sigma \vdash g_j\theta\}\big| \;\le\; \prod_i |\tau_i|,
\]
with equality exactly when no guard prunes. The bound is the program's,
not the grounder's: a room of $c$ cells whose adjacency is a predicate
rather than a table emits $32\,c^2$ events for four rules ($1.3$
million at $c = 200$), and no grounder can emit fewer. Any grounder must
at least read the facts its guards name and write its output,
$\Omega(|\text{facts read}| + |\mathrm{ground}(r)|)$. \emph{Guard-first
binding} reaches that bound for acyclic guard bodies: each variable's
domain is filtered by every fact-table guard that names it as soon as
it is bound, a semi-join per variable, so a prefix no row can complete is
never extended. The alternative that answers a guard at its last
variable is join-before-filter twice over --- the type product before the
guard, the whole fact list before the constants --- and costs $c^3$
against an output of $8c$. A guard body with a cycle (a triangle
$e(X,Y) \tensor e(Y,Z) \tensor e(Z,X)$) is bounded instead by the
fractional edge cover of its hypergraph (the AGM bound), and the
compiler reports it so that the cost is never a surprise.

\begin{center}\small
\begin{tabular}{lrrr}
\toprule
cells (static form) & 200 & 450 & 840 \\
\midrule
guard at its last variable & $22$ s & $248$ s & $29$ min \\
guard-first binding & $37$ ms & $94$ ms & $0.19$ s \\
\bottomrule
\end{tabular}
\end{center}

\paragraph{Emission and inference are EFD's object.} Once grounded,
the program's net exists, and from here the proof-net reading of
\cite{EFD26} applies verbatim: a site is a box, an occurrence edge is a
cut, the elimination order is a sequence of cut expansions, and exact
inference costs $O(n \cdot 2^{W})$ in the induced width $W$ of
\S\ref{sec:certs}. The two stages do not trade against each other.
Every room-shaped program, whatever its size, emits the same net at a
given horizon (the agent reaches the same cells), with a width far past
exact inference; a program written over \emph{roles} --- $k$ places the
cells are bound to --- has width $8$ at the same horizon. The width is
chosen by the program, as the event count is; what the compiler
chooses is only the order in which the program's own cost is paid.

Grounding cost is not inference cost, but the event count prices
everything downstream once: the load and memory of every runtime (about
a microsecond and a few hundred bytes per event), and the size of each
\textsc{choice} site's table in the net. A turn of the interpreter, by
contrast, is bounded by the candidates weighed at that state, and an
index by atom makes it independent of the program's size.

\subsection{The factored emission and the \texttt{reads} contract}
\label{sec:emission}

The emission (per micro-step): constants propagate --- control atoms
(chain and stage tokens) fold away, \emph{branch-independence itself}
being the criterion, with no control/state annotation anywhere. Steps
whose candidates or scopes touch branch-dependent atoms become
\textbf{sites}: categorical outcome variables whose CPT rows are the
normalized \textsc{choice} weights, with registry bindings evaluated on
synthesized snapshots restricted to their \textbf{declared scope} ---
the language-level contract
\texttt{weight f(C) reads \{rb(C), prev(D) if adj(C,D)\}}, guards
enumerated against the Horn world. The snapshot object \emph{raises} on
any out-of-scope read: the contract has teeth, and
Theorem~\ref{thm:trace} is why it must (an undeclared read breaks
footprint stability). Writes become SSA versions via deterministic
factors. Conservative static checks (multiplicities, schedule, double
production) reject at admission; constructs static analysis cannot
decide (stalled tiers, produce-onto-maybe-present) are emitted as
zero-mass rows and indicator guards, and the evaluation-time total-mass
check $Z = 1$ converts any \emph{reachable} violation into an error.
Conservative statics, exact runtime detection, never a silent
approximation.

\subsection{Certificates and the real-time doctrine}\label{sec:certs}

\texttt{certificate(query)} computes induced width and total factor-ops
from IR scopes alone --- no tables built, no binding called
(\stpinned{}: a counting registry observes zero calls at admission).
Budgets are platform-independent factor-ops, converted by one measured
constant per machine; over-budget
queries fall down the evaluator ladder (exact-support $\to$ factored
$\to$ fixed-$N$ sampling, whose cost is deterministic by construction).
Measured tracking on Game of Life instances: certified ops $10^4$ /
$5 \times 10^6$ / $7 \times 10^7$ against measured 0.00 / 0.04 / 0.40
s. The junction tree makes the elimination structure itself AOT: built
once per \texttt{program\_key}, two-pass calibration per evidence;
\texttt{unknowns} lift initial atoms to variables so \emph{one}
compiled structure serves every tick of a receding-horizon loop (the
bounded-slot-template deployment story), with per-tick evidence as
point factors.

\subsection{Honesty: what the cliff is and is not}\label{sec:cliff}

The exact filter dies where reachable support explodes --- synchronous
IPPC domains reach $2^n$ at $n \approx 14$--$18$. The factored back-end
restores exactness where induced width stays local (ring topologies,
bounded neighborhoods: $n = 24$ two-tick GoL in seconds), and
\emph{certifiably refuses} where it does not (the certificate is the
refusal). The sampler and THTS carry the rest with fixed budgets.

Measured on one machine, in the Python referee:

\begin{center}\small
\begin{tabular}{llll}
\toprule
evaluator & topology / instance & cost & note \\
\midrule
exact support filter & sync ring, 1-tick proj., $n=6$ & $<0.01$ s &
support $2^6$ \\
 & $n=12$ & 0.65 s & $2^{12}$; clean $\sim 4\times$ per $+2$ units \\
 & $n=14$ & 3.64 s & $2^{14}$ --- per-step planning tolerable \\
 & $n=16$ & 18.3 s & $2^{16}$ --- per-step planning dead \\
\addlinespace
factored / junction tree & sync ring, $n=16$, two ticks (96
micro-steps) & $<60$ s & exact past the filter's cliff;
sampler-refereed \\
 & GoL ring, $n=24$, two ticks & $<30$ s & exact; support filter would
need $2^{24}$ \\
 & GoL instance 5 ($n=16$), live play & 65 ms/decision &
certificate-priced, real-time compatible \\
\addlinespace
fixed-$N$ sampler & sync ring, $n=14$ & 1.1 s (abs err 0.008) & vs 3.8
s exact --- the crossover \\
 & sync ring, $n=20$ & 2.1 s & exact infeasible at $2^{20}$ \\
\bottomrule
\end{tabular}
\end{center}

The pattern the table quantifies: the exact filter's exponent is the
physics of support enumeration (topology-independent --- a ring at
$n = 16$ costs the same $2^{16}$ as a clique); the factored evaluator's
cost tracks \emph{induced width}, so it is topology-dependent --- local
neighborhoods (rings, grids) stay exact for horizons and sizes the
filter cannot touch, while densely-coupled domains blow the width and
are \emph{refused at admission} by the same certificate that admits the
local ones (certified ops vs measured seconds calibrate at one
constant, \S\ref{sec:certs}). The sampler's cost is deterministic by
construction and carries what both exact evaluators refuse. On IPPC
instances where exactness is in reach, the planner is provably optimal
(exact expectimax over the true measure: Navigation, all tested
instances; CrossingTraffic instance 1); Overcooked layouts reach expert
parity (185/185); nothing is claimed against PROST-class planners on
the large instances --- that comparison is \stopen{} and would be
sampling-based.

\subsection{C++ as IR interpreter; the compiler as verification
toolchain}\label{sec:cpp}

The deployment decision inverts the usual port: C++ never gains a
grounder, typechecker or min-fill. The compiled artifact (structure +
AOT schedules, serialized on the golden wire) is interpreted by a
$\sim$300-line dependency-free table-ops loop; CPT rows are filled at
run time by the registry bindings (they read the live world; that is
the one thing that must be native). Admission happened on the authoring
machine; the vehicle runs pre-certified artifacts with deterministic
cost. The authoritative compiler is \metisc{} (OCaml, standard library
only, one static binary, from lexing to the reference sampler); the
Python toolkit is the independent exact referee. Verification is a chain of custody: interpreter $\doteq$
filter (Thm~\ref{thm:preservation}--\ref{thm:filter} gates) $\doteq$
factored/junction tree ($10^{-9}$ parity, whole-tick and mid-cycle)
$\doteq$ exported artifact (independent brute-force reader) $\doteq$
C++ interpreter (golden parity on every corpus program, in both
toolchains) --- with pgmpy
belief propagation as an external referee on the exported net.

\subsection{Export: the EFD direction as a back-end}\label{sec:export}

\texttt{to\_uai} dumps the emitted net as a standard factor graph (UAI
MARKOV; pgmpy/libDAI/merlin all read it), evidence baked in. This is
the EFD direction --- program $\to$ net $\to$ BN --- obtained \emph{for
free} because the IR already is the net; it doubles as external
validation and as the data interface for parameter estimation
(\S\ref{sec:learning}).

\section{Learning: packs for structure, ports for parameters}
\label{sec:learning}

One invariant: learning never edits a rule in place. The catalog stays
the single reviewable, provenance-carrying artifact; every learned
artifact passes the same admission gates as a hand-authored one.

\paragraph{Structure (surprise $\to$ synthesis).} Zero predictive mass
is the structural signal \emph{by construction}: \textsc{choice}
weights multiply, so no parameter update can create missing support ---
a missing behavior is a missing rule, never a wrong number. Unmodeled
observations produce a small deterministic record
(\{\texttt{program\_key}, artifact versions, init multiset,
observation, $Z$\}); because evaluation is deterministic given these,
the offline side re-derives full context (support, event posteriors,
MAP explanation) --- the vehicle logs pointers, not traces. Candidates
(template enumeration over the typed vocabulary; the type system is the
hypothesis-space bound) enter as \texttt{provenance=surprise} packs
through containment admission.

\paragraph{Parameters (notching).} By Theorem~\ref{thm:trace}'s
factorization, each site is a multinomial over its declared scope ---
the emission literally exhibits the sufficient-statistic tables.
Observed traces give per-site counts (conjugate Dirichlet updates; the
hand prior and its strength are the pseudo-counts); partially observed
traces give expected counts via exact smoothing (event posteriors under
soft evidence) --- EM whose E-step is free. Hypothesis-level selection
over whole behavior rows is exact model filtering (posterior fed back
as next prior; planted-row identifiability $\ge 0.9$ recovery
\stpinned{}). Learned values ship as overlay packs with a
human-reviewable row sheet (\texttt{prior $\to$ posterior ($n$,
$\alpha_0$)} per scope assignment); the vehicle never learns in flight.
The batch interface to external estimators is (UAI net, traces) ---
data, not an RL arena; the simulators are trace generators for it.

\paragraph{The loop.} The loop that proposes packs --- a language model
drafting rules from surprise records, a gate admitting or refusing them
--- is the catalog paper's subject. Here two facts suffice: whatever
the model proposes passes the same admission as a hand-written rule,
and the world is driven by the admitted program, never by the model.
Learning changes the program; it never retrains anything.

\section{Mechanization}\label{sec:mech}

A Lean~4 development over mathlib ships in the public language
repository, directory \texttt{formal/}. This paper contains no
proof-assistant code: the development is the reference, and every
statement below names its Lean identifier. The development is
\texttt{sorry}-free, and the main theorems depend only on Lean's three
standard axioms (\texttt{propext}, \texttt{Classical.choice},
\texttt{Quot.sound}), checked with \texttt{\#print axioms}.
The design principle mirrors the compiler's: formalize the
\emph{content} of a theorem at the level of abstraction where its
mathematical difficulty lives, and leave the bookkeeping the
implementation already pins (parsers, hashing, wire formats) to the
golden gates.

\paragraph{Theorem~\ref{thm:image}, in full
(\texttt{Formal/Spn.lean}).} Ground programs are formalized as
name-carrying lists of clauses and links over ambient types of atoms,
stages and names; the SPN class as tuples whose transition components
are multisets of content records; the denotation forgets names and list
order; structural congruence is the existence of a renaming injective
on the program's names together with list permutations. The
development proves: \texttt{denote\_congr} (invariance: congruent
programs denote equal SPNs), \texttt{denote\_readBack} (surjectivity:
an explicit fresh-naming section, $\sem{U(N)} = N$ on the nose), and
\texttt{congr\_of\_denote\_eq} (injectivity up to congruence, for
well-formed programs --- distinct clause and link names); together
\texttt{denote\_eq\_iff\_congr}, the bijection of
Theorem~\ref{thm:image}. The delicate steps of the paper proof
reappear precisely where expected: the duplicate-content matching is a
permutation argument on lists with equal content multisets, and the
empty-stage and multiset-not-set remarks are forced by the statement
failing (as it should) under the alternative definitions.

\paragraph{Theorem~\ref{thm:preservation}, counting kernel
(\texttt{Formal/Counting.lean}).} An instance of clause $r$ at marking
$m$ is formalized per species: an injection of consume slots into
tokens and a function from read occurrences to tokens.
\texttt{card\_instance} proves the instance count equals
$\varphi(m,r)$ --- the falling factorial on consumed species (via
$\#(\alpha \hookrightarrow \beta) = \ff{|\beta|}{|\alpha|}$) times the
power on read species --- and \texttt{pushforward\_weight} proves the
per-clause pushforward identity $\sum_{\iota \in \mathrm{Inst}(r,m)}
w_r = w_r \cdot \varphi(m,r) = W_r(m)$, i.e.\ the interpreter's
instance-uniform draw and the DBN's \textsc{choice} weights induce the
same one-step clause distribution. This is the substantive combinatorial
content of Theorem~\ref{thm:preservation}; the surrounding product
induction is routine and remains on paper.

\paragraph{Theorem~\ref{thm:filter}
(\texttt{Formal/Filter.lean}).} Finitely-supported nonnegative
state distributions with a finite-branching stochastic kernel;
\texttt{filter\_eq\_marginal} proves by induction that the forward
filter computed over the (finite) support at each step equals the
$t$-step pushforward marginal --- restriction to the reachable support
loses no mass.

\paragraph{Theorem~\ref{thm:trace}, descent skeleton
(\texttt{Formal/Trace.lean}).} Events with an independence relation,
a deterministic update, and a local weight function; the two
hypotheses of the footprint-stability lemma --- independent events
commute as state transformers, and each event's weight factor is
unchanged by executing an independent event first --- are taken as the
definition of independence-compatibility. \texttt{weight\_exec\_swap}
proves the run weight (product of local factors along the run) and
final state are invariant under adjacent transposition of independent
events, and \texttt{traceEquiv\_weight\_eq} /
\texttt{traceEquiv\_exec\_eq} extend this to the Mazurkiewicz
equivalence closure: the measure descends to trace classes. The concrete instantiation --- that the factored emission's
sites and declared scopes satisfy the two hypotheses (the paper's
footprint-stability lemma) --- is exactly the \texttt{reads}-contract
argument of Theorem~\ref{thm:trace} and remains on paper, pinned
operationally by the emission's rejection gates.

\paragraph{Theorem~\ref{thm:mixture} and the cut bound
(\texttt{Formal/Belief.lean}).} Over the same finite state space and
rational masses as \texttt{Filter.lean}: \texttt{marginal\_sum} proves
the forward filter linear in its initial distribution, and
\texttt{hit\_mixture} that the instrument started from the world
mixture returns the mixture of the single-world answers for any kernel,
outcome set and horizon --- with the absorbing kernel \texttt{absorb},
the first-passage number of \S\ref{sec:instrument}.
\texttt{hit\_cut\_bounds} proves the interval of \S\ref{sec:instrument}:
a filter that sets states aside at each step returns $\mathit{lo} \le
\text{exact} \le \mathit{lo} + \mathit{CUT}$ under a nonnegative
stochastic kernel, through two invariants of the dropped part of the
distribution --- it stays nonnegative (\texttt{dropped\_nonneg}) and
its total is exactly the mass set aside (\texttt{dropped\_total}).
Proposition~\ref{prop:cells} (multilinearity) remains on paper.

\paragraph{What ``automated'' means.} A belief table's side
conditions --- every group's masses nonnegative and summing to one, a
level port's cut points strictly increasing in $(0,1)$ --- are a
decidable predicate \texttt{Sheet.ok} on concrete data, closed by
normalization (\texttt{norm\_num}) in the development's examples. That
is the form the compiler can emit per package, next to the goldens: a
statement per table, checked by the kernel, with
Theorem~\ref{thm:mixture} proved once above them. What cannot be
emitted this way is a port as a function of the step
(\S\ref{sec:bystep}): the language refuses it, so there is nothing to
check.

\paragraph{What is deliberately not mechanized.} Propositions
\ref{prop:quiescence} and \ref{prop:bpn} are statements about the
emission as implemented; their assurance story is the admission-time
checker that runs at every serialization (\stpinned{}), which checks
the actual artifact rather than a model of it. Mechanizing the checker
itself (verify the verifier) and Conjecture~\ref{conj:typing}'s
compositional typing are the natural next targets
(\S\ref{sec:open}).

\section{Related work and provenance}\label{sec:related}

\paragraph{The direct anchor.} EFD \cite{EFD24,EFD26}: BNs
$\leftrightarrow$ $\mllb$ proof-nets, inference = net evaluation in
probabilistic coherence spaces, cost-matched; \cite{DEF26} adds
compositional typing (quantum and classical). Acclavio--Maieli
\cite{AM26}: proof \emph{search} in probabilistic LL computing BN
inference --- the operational face. Our audit trail is explicit: the
$\lolli$-as-V-joint observation predates our contact with this
literature but was never published by us; priority on the static
correspondence is theirs. What is ours is the dynamic, staged,
executable direction (\S\ref{sec:intro} items 1--9) and its
compilation.

\paragraph{Foundations used.} LL$\leftrightarrow$Petri:
\cite{GG89,EW93}; MSR/Ceptre lineage: \cite{CS09,Mar15,CLF};
probabilistic coherence spaces and weighted relational models:
\cite{Gir04,DE11,EPT18,LMMP13}; Markov categories: \cite{Fri20,CJ19};
true-concurrency probability: \cite{AB06,Win13,CCPW18,Paq21};
GSPN/mass action: \cite{Gil77,GSPN}; knowledge compilation: ProbLog
$\to$ d-DNNF/SDD \cite{FdB15,Dar11}; classical-logic contrast: Markov
Logic Networks \cite{RD06} ground FOL into Markov nets --- the linear
analogue sat unoccupied until the EFD line; planning-as-inference:
\cite{TS06,Lev18}; mechanization: Lean 4 \cite{MU21} and mathlib
\cite{mathlib20}. Behavior-tree engineering contrast: BT.cpp gives
industry the component model but no measure; our framework keeps the
component seams (registry, blackboard, ports) and adds semantics, and
its interface packages are the contract half of that component model.

\paragraph{The belief side.} A program with belief ports is a
finite-horizon belief MDP \cite{KLC98} evaluated at one tick over a
finite quantisation. What is particular here is that the quantisation is
the program's \emph{input type} --- the level atom --- so the belief MDP
is certified by a finite sheet over cells rather than by a bound over a
continuous belief space, and the filter that moves the belief stays
outside the certified object.

\section{Open problems}\label{sec:open}

\begin{enumerate}[leftmargin=2em]
\item The full trajectory-level statement of
  Theorem~\ref{thm:preservation} (the product induction with
  normalization over enabled sets) and the concrete site/scope
  instantiation of Theorem~\ref{thm:trace} over a formalized emission.
\item Branching-cell measure beyond $\cclass$: import \cite{AB06} to
  define trace-invariant weights for programs the emission rejects;
  does the general staged language admit a canonical net measure?
\item Conjecture~\ref{conj:typing}: mechanize the \cite{DEF26}
  well-behaved-composition typing derivations over the emitted
  artifact, and verify the verifier: re-prove the artifact checker of
  Proposition~\ref{prop:bpn} in Lean against the formalized IR.
\item Proof theory of quiescence in the general class
  (subexponentials, fixpoints, or inhibitor-net semantics).
\item Ports as functions (\S\ref{sec:belief}): the instrument as a
  filter over a hidden chain, updated by what each trace observes. What
  replaces multilinearity, and over which parameters a certificate can
  then be stated.
\item The believed position of an unseen agent as a hazard: a
  predicted position counts as a hazard in a safety record and should
  not count as one in a verdict; a belief-aware form of the record.
\item Faithfulness: which queries determine a program in $\cclass$?
  And the engineering comparisons --- sampling-based planners on large
  benchmark instances, complementary-pair merging and lifted emission
  --- which this paper does not claim.
\end{enumerate}

\begin{thebibliography}{99}\small

\bibitem[AB06]{AB06} S. Abbes, A. Benveniste. True-concurrency
probabilistic models: branching cells and distributed probabilities.
2006/08.

\bibitem[AM26]{AM26} M. Acclavio, R. Maieli. Probabilistic Linear
Logic Programming with an Application to Bayesian Network
Computations. arXiv:2601.13270.

\bibitem[CCPW18]{CCPW18} S. Castellan, P. Clairambault, H. Paquet, G.
Winskel. The concurrent game semantics of probabilistic PCF. LICS
2018.

\bibitem[CJ19]{CJ19} K. Cho, B. Jacobs. Disintegration and Bayesian
inversion via string diagrams. MSCS 2019.

\bibitem[CLF]{CLF} K. Watkins, I. Cervesato, F. Pfenning, D. Walker. A
concurrent logical framework. CMU-CS-02-101.

\bibitem[CS09]{CS09} I. Cervesato, A. Scedrov. Relating state-based
and process-based concurrency through linear logic. Inf. \& Comp.
2009.

\bibitem[Dar11]{Dar11} A. Darwiche. SDD: a new canonical
representation of propositional knowledge bases. IJCAI 2011.

\bibitem[DE11]{DE11} V. Danos, T. Ehrhard. Probabilistic coherence
spaces as a model of higher-order probabilistic computation. Inf. \&
Comp. 209(6), 2011.

\bibitem[DEF26]{DEF26} R. Di Guardia, T. Ehrhard, C. Faggian. Quantum
Bayesian Networks: Compositionality and Typing via Linear Logic.
arXiv:2604.26059.

\bibitem[EFD24]{EFD24} T. Ehrhard, C. Faggian, R. Di Guardia. Bayesian
Networks and Proof-Nets: a Proof-Theoretical Account of Bayesian
Inference. arXiv:2412.20540.

\bibitem[EFD26]{EFD26} R. Di Guardia, T. Ehrhard, A. Evrard, C.
Faggian. Bayesian Networks and Proof-Nets: the Proof-Theory of
Bayesian Inference. arXiv:2602.04045.

\bibitem[EPT18]{EPT18} T. Ehrhard, M. Pagani, C. Tasson. Full
abstraction for probabilistic PCF. J. ACM 2018.

\bibitem[EW93]{EW93} U. Engberg, G. Winskel. Linear logic on Petri
nets. 1990/93.

\bibitem[FdB+15]{FdB15} D. Fierens, G. Van den Broeck, et al.
Inference and learning in probabilistic logic programs using weighted
boolean formulas. TPLP 2015.

\bibitem[Fri20]{Fri20} T. Fritz. A synthetic approach to Markov
kernels, conditional independence and theorems on sufficient
statistics. Adv. Math. 2020.

\bibitem[GG89]{GG89} C. Gunter, V. Gehlot. Nets as tensor theories.
1989.

\bibitem[Gil77]{Gil77} D. Gillespie. Exact stochastic simulation of
coupled chemical reactions. J. Phys. Chem. 1977.

\bibitem[Gir87]{Gir87} J.-Y. Girard. Linear logic. Theoretical
Computer Science 50(1), 1987.

\bibitem[Gir04]{Gir04} J.-Y. Girard. Between logic and quantic: a
tract. 2004.

\bibitem[GSPN]{GSPN} M. Ajmone Marsan et al. Modelling with
Generalized Stochastic Petri Nets. Wiley 1995.

\bibitem[KLC98]{KLC98} L. P. Kaelbling, M. L. Littman, A. R.
Cassandra. Planning and acting in partially observable stochastic
domains. Artificial Intelligence 101, 1998.

\bibitem[Kol33]{Kol33} A. N. Kolmogorov. Grundbegriffe der
Wahrscheinlichkeitsrechnung. Springer 1933.

\bibitem[Lev18]{Lev18} S. Levine. Reinforcement learning and control
as probabilistic inference. arXiv:1805.00909.

\bibitem[LMMP13]{LMMP13} J. Laird, G. Manzonetto, G. McCusker, M.
Pagani. Weighted relational models of typed lambda-calculi. LICS 2013.

\bibitem[Mar15]{Mar15} C. Martens. Ceptre: a language for modeling
generative interactive systems. AIIDE 2015.

\bibitem[mathlib20]{mathlib20} The mathlib Community. The Lean
mathematical library. CPP 2020.

\bibitem[MU21]{MU21} L. de Moura, S. Ullrich. The Lean 4 theorem
prover and programming language. CADE 2021.

\bibitem[Paq21]{Paq21} H. Paquet. Bayesian strategies: probabilistic
programs as generalised graphical models. ESOP 2021.

\bibitem[RD06]{RD06} M. Richardson, P. Domingos. Markov logic
networks. MLJ 2006.

\bibitem[TS06]{TS06} M. Toussaint, A. Storkey. Probabilistic inference
for solving discrete and continuous state Markov decision processes.
ICML 2006.

\bibitem[Win13]{Win13} G. Winskel. Distributed probabilistic and
quantum strategies. MFPS 2013.

\end{thebibliography}

\end{document}
